\section{Characters and orthogonality relations}
Let now
\[
    k=\mathbb{C}.
\]
Let \((V,\rho)\) be a finite-dimensional complex representation of \(G\), so
\[
    \rho:G\longrightarrow \operatorname{GL}(V).
\]
After choosing a basis of \(V\), this may be written as
\[
    \rho:G\longrightarrow \operatorname{GL}_d(\mathbb{C}),
    \qquad
    d=\dim_{\mathbb{C}}V.
\]

\begin{definition}[Character]
    Given a complex representation \((V,\rho)\) of \(G\), its character is the
    function
    \[
        \chi_V:G\longrightarrow \mathbb{C}
    \]
    defined by
    \[
        \chi_V(g)=\operatorname{tr}\rho(g).
    \]
\end{definition}

\begin{proposition}
    Let \(V\) and \(W\) be finite-dimensional complex representations of \(G\).
    Then:

    \begin{enumerate}
        \item \(\chi_V\) is a class function.

        \item For every \(g\in G\),
              \[
                  \chi_V(g^{-1})=\overline{\chi_V(g)}.
              \]

        \item If \(V\cong W\) as \(\mathbb{C}G\)-modules, then
              \[
                  \chi_V=\chi_W.
              \]

        \item We have
              \[
                  \chi_{V\oplus W}=\chi_V+\chi_W,
                  \qquad
                  \chi_{V\otimes W}=\chi_V\chi_W.
              \]

        \item For the dual representation \(V^*\),
              \[
                  \chi_{V^*}(g)
                  =
                  \chi_V(g^{-1})
                  =
                  \overline{\chi_V(g)}.
              \]
              Hence
              \[
                  \chi_{V^*}=\overline{\chi_V}.
              \]

        \item For the representation \(\operatorname{Hom}_{\mathbb{C}}(V,W)\),
              \[
                  \chi_{\operatorname{Hom}_{\mathbb{C}}(V,W)}(g)
                  =
                  \chi_V(g^{-1})\chi_W(g).
              \]
              Equivalently,
              \[
                  \chi_{\operatorname{Hom}_{\mathbb{C}}(V,W)}
                  =
                  \overline{\chi_V}\chi_W.
              \]

        \item One has
              \[
                  \chi_V(1)=\dim_{\mathbb{C}}V.
              \]
              Moreover,
              \[
                  |\chi_V(g)|\leq \dim_{\mathbb{C}}V.
              \]
              Equality holds if and only if \(g\) acts on \(V\) as a scalar. In particular,
              \[
                  \chi_V(g)=\dim_{\mathbb{C}}V
              \]
              if and only if \(g\) acts trivially on \(V\).
    \end{enumerate}
\end{proposition}

\begin{proof}
    For \((1)\), let \(h,g\in G\). Then
    \[
        \rho(hgh^{-1})
        =
        \rho(h)\rho(g)\rho(h)^{-1}.
    \]
    Since trace is invariant under conjugation,
    \[
        \chi_V(hgh^{-1})
        =
        \operatorname{tr}\rho(hgh^{-1})
        =
        \operatorname{tr}\rho(g)
        =
        \chi_V(g).
    \]
    Thus \(\chi_V\) is a class function.

    For \((2)\), since \(g\) has finite order, \(\rho(g)\) also has finite order.
    Hence its eigenvalues are roots of unity. If the eigenvalues of \(\rho(g)\)
    are
    \[
        \lambda_1,\dots,\lambda_d,
    \]
    then the eigenvalues of \(\rho(g^{-1})=\rho(g)^{-1}\) are
    \[
        \lambda_1^{-1},\dots,\lambda_d^{-1}.
    \]
    Since each \(\lambda_i\) is a root of unity, we have
    \[
        \lambda_i^{-1}=\overline{\lambda_i}.
    \]
    Therefore
    \[
        \chi_V(g^{-1})
        =
        \sum_{i=1}^d\lambda_i^{-1}
        =
        \sum_{i=1}^d\overline{\lambda_i}
        =
        \overline{\chi_V(g)}.
    \]

    For \((3)\), if \(V\cong W\), then after choosing compatible bases, the matrices
    of \(\rho_V(g)\) and \(\rho_W(g)\) are conjugate for every \(g\in G\). Hence
    their traces are equal.

    For \((4)\), the matrix of \(g\) on \(V\oplus W\) is block diagonal:
    \[
        \rho_{V\oplus W}(g)
        =
        \begin{pmatrix}
            \rho_V(g) & 0         \\
            0         & \rho_W(g)
        \end{pmatrix}.
    \]
    Therefore
    \[
        \chi_{V\oplus W}(g)
        =
        \chi_V(g)+\chi_W(g).
    \]
    Also,
    \[
        \rho_{V\otimes W}(g)
        =
        \rho_V(g)\otimes \rho_W(g),
    \]
    so
    \[
        \operatorname{tr}(\rho_V(g)\otimes \rho_W(g))
        =
        \operatorname{tr}\rho_V(g)\operatorname{tr}\rho_W(g).
    \]
    Thus
    \[
        \chi_{V\otimes W}(g)=\chi_V(g)\chi_W(g).
    \]

    For \((5)\), the action of \(g\) on \(V^*\) is dual to the action of \(g^{-1}\)
    on \(V\). Hence
    \[
        \chi_{V^*}(g)
        =
        \operatorname{tr}\rho_V(g^{-1})
        =
        \chi_V(g^{-1})
        =
        \overline{\chi_V(g)}.
    \]

    For \((6)\), we use the natural identification
    \[
        \operatorname{Hom}_{\mathbb{C}}(V,W)
        \cong
        V^*\otimes W.
    \]
    Therefore, by \((4)\) and \((5)\),
    \[
        \chi_{\operatorname{Hom}_{\mathbb{C}}(V,W)}
        =
        \chi_{V^*}\chi_W
        =
        \overline{\chi_V}\chi_W.
    \]
    Equivalently, for every \(g\in G\),
    \[
        \chi_{\operatorname{Hom}_{\mathbb{C}}(V,W)}(g)
        =
        \chi_V(g^{-1})\chi_W(g).
    \]

    Finally, for \((7)\), we have
    \[
        \rho(1)=\operatorname{id}_V,
    \]
    so
    \[
        \chi_V(1)=\operatorname{tr}(\operatorname{id}_V)=\dim_{\mathbb{C}}V.
    \]
    Let the eigenvalues of \(\rho(g)\) be
    \[
        \lambda_1,\dots,\lambda_d.
    \]
    Since \(g\) has finite order, each \(\lambda_i\) is a root of unity, hence
    \[
        |\lambda_i|=1.
    \]
    Therefore
    \[
        |\chi_V(g)|
        =
        \left|\lambda_1+\cdots+\lambda_d\right|
        \leq
        |\lambda_1|+\cdots+|\lambda_d|
        =
        d.
    \]
    Equality holds if and only if all \(\lambda_i\) have the same argument, that is,
    if and only if all \(\lambda_i\) are equal. Since \(\rho(g)\) is diagonalizable,
    this is equivalent to saying that \(\rho(g)\) is a scalar operator.

    Moreover,
    \[
        \chi_V(g)=d
    \]
    if and only if all eigenvalues of \(\rho(g)\) are equal to \(1\), which is
    equivalent to
    \[
        \rho(g)=\operatorname{id}_V.
    \]
    Thus \(g\) acts trivially on \(V\).
\end{proof}

For simplicity, write
\[
    \operatorname{CF}(G)
    =
    \operatorname{CF}(G,\mathbb{C}).
\]
Thus \(\operatorname{CF}(G)\) is the \(\mathbb{C}\)-vector space of complex
class functions on \(G\).

\begin{definition}
    For \(f_1,f_2\in \operatorname{CF}(G)\), define
    \[
        \langle f_1,f_2\rangle
        =
        \frac{1}{|G|}
        \sum_{g\in G}
        f_1(g)\overline{f_2(g)}.
    \]
\end{definition}

Since \(f_1\) and \(f_2\) are class functions, the above sum may be written as
a sum over conjugacy classes. Let \(\operatorname{Cl}(G)\) denote the set of
conjugacy classes of \(G\). For each \(C\in \operatorname{Cl}(G)\), choose a
representative \(g_C\in C\). Then
\[
    \begin{aligned}
        \langle f_1,f_2\rangle
         & =
        \frac{1}{|G|}
        \sum_{C\in \operatorname{Cl}(G)}
        |C| f_1(g_C)\overline{f_2(g_C)} \\
         & =
        \frac{1}{|G|}
        \sum_{C\in \operatorname{Cl}(G)}
        \frac{|G|}{|C_G(g_C)|}
        f_1(g_C)\overline{f_2(g_C)}     \\
         & =
        \sum_{C\in \operatorname{Cl}(G)}
        \frac{1}{|C_G(g_C)|}
        f_1(g_C)\overline{f_2(g_C)}.
    \end{aligned}
\]
Here
\[
    C_G(g_C)
    =
    \{\,h\in G\mid hg_C=g_Ch\,\}
\]
is the centralizer of \(g_C\) in \(G\), and we used the orbit-stabilizer formula
\[
    |C|
    =
    \frac{|G|}{|C_G(g_C)|}.
\]

\begin{proposition}
    The pairing
    \[
        \langle -,-\rangle:
        \operatorname{CF}(G)\times \operatorname{CF}(G)
        \longrightarrow \mathbb{C}
    \]
    defined by
    \[
        \langle f_1,f_2\rangle
        =
        \frac{1}{|G|}
        \sum_{g\in G} f_1(g)\overline{f_2(g)}
    \]
    is a Hermitian inner product on the complex vector space
    \(\operatorname{CF}(G)\).

    Moreover, if \(f_1\) and \(f_2\) are characters of complex representations of
    \(G\), then
    \[
        \langle f_1,f_2\rangle\in \mathbb{R}.
    \]
\end{proposition}

\begin{proof}
    It is clear that the pairing is linear in the first variable and conjugate
    linear in the second variable. Also,
    \[
        \begin{aligned}
            \overline{\langle f_1,f_2\rangle}
             & =
            \overline{
                \frac{1}{|G|}\sum_{g\in G} f_1(g)\overline{f_2(g)}
            }                                    \\
             & =
            \frac{1}{|G|}
            \sum_{g\in G}\overline{f_1(g)}f_2(g) \\
             & =
            \langle f_2,f_1\rangle.
        \end{aligned}
    \]
    Finally,
    \[
        \langle f,f\rangle
        =
        \frac{1}{|G|}
        \sum_{g\in G}|f(g)|^2\geq 0,
    \]
    and equality holds if and only if \(f(g)=0\) for all \(g\in G\), that is,
    \(f=0\). Thus \(\langle -,-\rangle\) is a Hermitian inner product.

    Now suppose that \(f_1\) and \(f_2\) are characters. For every \(g\in G\), one
    has
    \[
        \overline{f_i(g)}=f_i(g^{-1}),
        \qquad i=1,2.
    \]
    Hence
    \[
        \begin{aligned}
            \langle f_1,f_2\rangle
             & =
            \frac{1}{|G|}
            \sum_{g\in G} f_1(g)\overline{f_2(g)} \\
             & =
            \frac{1}{|G|}
            \sum_{g\in G} f_1(g)f_2(g^{-1})       \\
             & =
            \frac{1}{|G|}
            \sum_{g\in G} f_1(g^{-1})f_2(g)       \\
             & =
            \frac{1}{|G|}
            \sum_{g\in G} \overline{f_1(g)}f_2(g) \\
             & =
            \langle f_2,f_1\rangle.
        \end{aligned}
    \]
    But Hermitian symmetry gives
    \[
        \langle f_2,f_1\rangle
        =
        \overline{\langle f_1,f_2\rangle}.
    \]
    Therefore
    \[
        \langle f_1,f_2\rangle
        =
        \overline{\langle f_1,f_2\rangle},
    \]
    and hence
    \[
        \langle f_1,f_2\rangle\in \mathbb{R}.
    \]
\end{proof}

\begin{lemma}
    Let \(V\) be a finite-dimensional \(\mathbb{C}G\)-module. Define
    \[
        V^G
        =
        \{\,v\in V\mid gv=v \text{ for all } g\in G\,\}.
    \]
    Then \(V^G\) is a \(\mathbb{C}G\)-submodule of \(V\), and
    \[
        \dim_{\mathbb{C}} V^G
        =
        \langle \chi_V,1_G\rangle
        =
        \frac{1}{|G|}
        \sum_{g\in G}\chi_V(g),
    \]
    where \(1_G\) denotes the character of the trivial representation.
\end{lemma}

\begin{proof}
    It is clear that \(V^G\) is a complex subspace of \(V\). Moreover, if
    \(v\in V^G\), then for every \(h\in G\),
    \[
        hv=v.
    \]
    Thus \(V^G\) is stable under the action of \(G\), and the induced action on
    \(V^G\) is trivial. Hence \(V^G\) is a \(\mathbb{C}G\)-submodule of \(V\).

    Define a linear map
    \[
        \pi:V\longrightarrow V
    \]
    by
    \[
        \pi(v)
        =
        \left(
        \frac{1}{|G|}
        \sum_{g\in G}g
        \right)v
        =
        \frac{1}{|G|}
        \sum_{g\in G}gv.
    \]
    We claim that
    \[
        \operatorname{Im}\pi=V^G.
    \]

    First, for \(h\in G\), we have
    \[
        \begin{aligned}
            h\pi(v)
             & =
            h\left(
            \frac{1}{|G|}
            \sum_{g\in G}gv
            \right)          \\
             & =
            \frac{1}{|G|}
            \sum_{g\in G}hgv \\
             & =
            \frac{1}{|G|}
            \sum_{g\in G}gv  \\
             & =
            \pi(v),
        \end{aligned}
    \]
    because \(g\mapsto hg\) is a permutation of \(G\). Hence
    \[
        \operatorname{Im}\pi\subseteq V^G.
    \]

    Conversely, if \(v\in V^G\), then
    \[
        \pi(v)
        =
        \frac{1}{|G|}
        \sum_{g\in G}gv
        =
        \frac{1}{|G|}
        \sum_{g\in G}v
        =
        v.
    \]
    Thus \(V^G\subseteq \operatorname{Im}\pi\). Therefore
    \[
        \operatorname{Im}\pi=V^G.
    \]

    Also,
    \[
        \pi^2
        =
        \left(
        \frac{1}{|G|}
        \sum_{g\in G}g
        \right)
        \left(
        \frac{1}{|G|}
        \sum_{h\in G}h
        \right)
        =
        \frac{1}{|G|^2}
        \sum_{g,h\in G}gh.
    \]
    For each \(x\in G\), there are exactly \(|G|\) pairs \((g,h)\in G\times G\)
    such that \(gh=x\). Hence
    \[
        \pi^2
        =
        \frac{1}{|G|}
        \sum_{x\in G}x
        =
        \pi.
    \]
    Thus \(\pi\) is a projection onto \(V^G\). Therefore
    \[
        V=\ker\pi\oplus \operatorname{Im}\pi,
        \qquad
        \pi|_{\operatorname{Im}\pi}=\operatorname{id}_{\operatorname{Im}\pi}.
    \]
    It follows that
    \[
        \operatorname{tr}(\pi)=\dim_{\mathbb{C}}\operatorname{Im}\pi
        =
        \dim_{\mathbb{C}}V^G.
    \]

    On the other hand,
    \[
        \begin{aligned}
            \operatorname{tr}(\pi)
             & =
            \operatorname{tr}\left(
            \frac{1}{|G|}
            \sum_{g\in G}\rho(g)
            \right)                                 \\
             & =
            \frac{1}{|G|}
            \sum_{g\in G}\operatorname{tr}(\rho(g)) \\
             & =
            \frac{1}{|G|}
            \sum_{g\in G}\chi_V(g).
        \end{aligned}
    \]
    Since \(1_G(g)=1\) for all \(g\in G\), we have
    \[
        \frac{1}{|G|}
        \sum_{g\in G}\chi_V(g)
        =
        \langle \chi_V,1_G\rangle.
    \]
    Therefore
    \[
        \dim_{\mathbb{C}}V^G
        =
        \langle \chi_V,1_G\rangle.
    \]
\end{proof}

\begin{lemma}
    Let \(k\) be any field, and let \(V,W\) be \(kG\)-modules. Give
    \(\operatorname{Hom}_k(V,W)\) the \(G\)-action defined by
    \[
        (g\varphi)(v)
        =
        g\varphi(g^{-1}v),
        \qquad
        g\in G,\ \varphi\in\operatorname{Hom}_k(V,W),\ v\in V.
    \]
    Then
    \[
        \operatorname{Hom}_k(V,W)^G
        =
        \operatorname{Hom}_{kG}(V,W).
    \]
\end{lemma}

\begin{proof}
    Let \(\varphi\in \operatorname{Hom}_k(V,W)\). Then
    \[
        \varphi\in \operatorname{Hom}_k(V,W)^G
    \]
    if and only if
    \[
        g\varphi=\varphi
    \]
    for all \(g\in G\). By definition of the action, this means
    \[
        g\varphi(g^{-1}v)=\varphi(v)
    \]
    for all \(g\in G\) and \(v\in V\). Replacing \(v\) by \(gv\), this is
    equivalent to
    \[
        g\varphi(v)=\varphi(gv)
    \]
    for all \(g\in G\) and \(v\in V\). This is precisely the condition that
    \(\varphi\) be a \(kG\)-module homomorphism. Hence
    \[
        \operatorname{Hom}_k(V,W)^G
        =
        \operatorname{Hom}_{kG}(V,W).
    \]
\end{proof}

\begin{theorem}[Row orthogonality]
    Let \(V\) and \(W\) be simple \(\mathbb{C}G\)-modules. Then
    \[
        \langle \chi_V,\chi_W\rangle
        =
        \begin{cases}
            1, & V\cong W,     \\
            0, & V\not\cong W.
        \end{cases}
    \]
\end{theorem}

\begin{proof}
    By Schur's lemma,
    \[
        \dim_{\mathbb{C}}\operatorname{Hom}_{\mathbb{C}G}(V,W)
        =
        \begin{cases}
            1, & V\cong W,     \\
            0, & V\not\cong W.
        \end{cases}
    \]
    On the other hand, by the previous lemma,
    \[
        \operatorname{Hom}_{\mathbb{C}G}(V,W)
        =
        \operatorname{Hom}_{\mathbb{C}}(V,W)^G.
    \]
    Therefore
    \[
        \begin{aligned}
            \dim_{\mathbb{C}}\operatorname{Hom}_{\mathbb{C}G}(V,W)
             & =
            \dim_{\mathbb{C}}\operatorname{Hom}_{\mathbb{C}}(V,W)^G \\
             & =
            \left\langle
            \chi_{\operatorname{Hom}_{\mathbb{C}}(V,W)},
            1_G
            \right\rangle.
        \end{aligned}
    \]
    But
    \[
        \operatorname{Hom}_{\mathbb{C}}(V,W)
        \cong
        V^*\otimes W
    \]
    as \(\mathbb{C}G\)-modules. Hence
    \[
        \chi_{\operatorname{Hom}_{\mathbb{C}}(V,W)}
        =
        \chi_{V^*}\chi_W
        =
        \overline{\chi_V}\chi_W.
    \]
    Thus
    \[
        \begin{aligned}
            \dim_{\mathbb{C}}\operatorname{Hom}_{\mathbb{C}G}(V,W)
             & =
            \frac{1}{|G|}
            \sum_{g\in G}\overline{\chi_V(g)}\chi_W(g) \\
             & =
            \langle \chi_W,\chi_V\rangle.
        \end{aligned}
    \]
    Since \(\chi_V\) and \(\chi_W\) are characters, the previous proposition gives
    \[
        \langle \chi_W,\chi_V\rangle
        =
        \langle \chi_V,\chi_W\rangle.
    \]
    Therefore
    \[
        \langle \chi_V,\chi_W\rangle
        =
        \dim_{\mathbb{C}}\operatorname{Hom}_{\mathbb{C}G}(V,W).
    \]
    By Schur's lemma, this dimension is \(1\) if \(V\cong W\), and \(0\) otherwise.
    Hence
    \[
        \langle \chi_V,\chi_W\rangle
        =
        \begin{cases}
            1, & V\cong W,     \\
            0, & V\not\cong W.
        \end{cases}
    \]
\end{proof}

\begin{definition}
    The set of irreducible characters of \(G\) is denoted by
    \[
        \operatorname{Irr}(G).
    \]
    Thus
    \[
        \operatorname{Irr}(G)
        =
        \{\,\chi_S\mid S \text{ is a simple } \mathbb{C}G\text{-module}\,\}.
    \]
    Equivalently, \(\operatorname{Irr}(G)\) is the set of characters afforded by
    simple complex representations of \(G\).
\end{definition}

\begin{corollary}
    The set \(\operatorname{Irr}(G)\) is \(\mathbb{C}\)-linearly independent in
    \(\operatorname{CF}(G)\).
\end{corollary}

\begin{proof}
    Suppose that
    \[
        \sum_{\chi\in \operatorname{Irr}(G)} a_\chi \chi=0,
        \qquad a_\chi\in \mathbb{C}.
    \]
    Let \(\psi\in \operatorname{Irr}(G)\). Taking the inner product with \(\psi\),
    and using the row orthogonality relations, we get
    \[
        0
        =
        \left\langle
        \sum_{\chi\in \operatorname{Irr}(G)} a_\chi \chi,\psi
        \right\rangle
        =
        \sum_{\chi\in \operatorname{Irr}(G)} a_\chi
        \langle \chi,\psi\rangle
        =
        a_\chi.
    \]
    Since this holds for every \(\psi\in \operatorname{Irr}(G)\), all coefficients
    are zero. Hence \(\operatorname{Irr}(G)\) is \(\mathbb{C}\)-linearly
    independent.
\end{proof}

\begin{corollary}
    The set \(\operatorname{Irr}(G)\) is a \(\mathbb{C}\)-basis of
    \(\operatorname{CF}(G)\).
\end{corollary}

\begin{proof}
    Let
    \[
        \operatorname{Irr}(G)=\{\chi_1,\dots,\chi_r\},
    \]
    where \(\chi_i\) is afforded by the simple \(\mathbb{C}G\)-module \(S_i\).
    Since \(\mathbb{C}\) is a splitting field for \(G\), the Artin--Wedderburn
    decomposition gives
    \[
        \mathbb{C}G
        \cong
        M_{d_1}(\mathbb{C})
        \oplus
        \cdots
        \oplus
        M_{d_r}(\mathbb{C}),
    \]
    where
    \[
        d_i=\dim_{\mathbb{C}}S_i.
    \]
    Taking centers, we get
    \[
        Z(\mathbb{C}G)
        \cong
        Z(M_{d_1}(\mathbb{C}))
        \oplus
        \cdots
        \oplus
        Z(M_{d_r}(\mathbb{C})).
    \]
    Since
    \[
        Z(M_{d_i}(\mathbb{C}))=\mathbb{C}I_{d_i},
    \]
    we have
    \[
        \dim_{\mathbb{C}}Z(\mathbb{C}G)=r.
    \]
    On the other hand, we have already seen that
    \[
        \operatorname{CF}(G)\cong Z(\mathbb{C}G)
    \]
    as \(\mathbb{C}\)-vector spaces. Therefore
    \[
        \dim_{\mathbb{C}}\operatorname{CF}(G)
        =
        r
        =
        |\operatorname{Irr}(G)|.
    \]
    Since \(\operatorname{Irr}(G)\) is linearly independent and has cardinality
    equal to \(\dim_{\mathbb{C}}\operatorname{CF}(G)\), it is a basis of
    \(\operatorname{CF}(G)\).
\end{proof}

\begin{corollary}
    The number of irreducible complex characters of \(G\) is equal to the number
    of conjugacy classes of \(G\). That is,
    \[
        |\operatorname{Irr}(G)|
        =
        |\operatorname{Cl}(G)|.
    \]
\end{corollary}

\begin{proof}
    Since \(\operatorname{Irr}(G)\) is a basis of \(\operatorname{CF}(G)\), we have
    \[
        |\operatorname{Irr}(G)|
        =
        \dim_{\mathbb{C}}\operatorname{CF}(G).
    \]
    But \(\operatorname{CF}(G)\) has a basis given by the characteristic functions
    of the conjugacy classes. Hence
    \[
        \dim_{\mathbb{C}}\operatorname{CF}(G)
        =
        |\operatorname{Cl}(G)|.
    \]
    Therefore
    \[
        |\operatorname{Irr}(G)|
        =
        |\operatorname{Cl}(G)|.
    \]
\end{proof}

\begin{lemma}
    One has
    \[
        |G|
        =
        \sum_{\chi\in \operatorname{Irr}(G)} \chi(1)^2.
    \]
\end{lemma}

\begin{proof}
    Let
    \[
        S_1,\dots,S_r
    \]
    be a complete set of pairwise non-isomorphic simple \(\mathbb{C}G\)-modules,
    and let
    \[
        \chi_i=\chi_{S_i},
        \qquad
        d_i=\dim_{\mathbb{C}}S_i.
    \]
    As a left \(\mathbb{C}G\)-module, the regular module decomposes as
    \[
        \mathbb{C}G
        \cong
        S_1^{\oplus d_1}
        \oplus
        \cdots
        \oplus
        S_r^{\oplus d_r}.
    \]
    Taking dimensions gives
    \[
        |G|
        =
        \dim_{\mathbb{C}}\mathbb{C}G
        =
        \sum_{i=1}^r d_i\dim_{\mathbb{C}}S_i
        =
        \sum_{i=1}^r d_i^2.
    \]
    Since
    \[
        d_i=\chi_i(1),
    \]
    we obtain
    \[
        |G|
        =
        \sum_{i=1}^r \chi_i(1)^2
        =
        \sum_{\chi\in \operatorname{Irr}(G)} \chi(1)^2.
    \]
\end{proof}

\begin{corollary}
    Let \(V\) be a finite-dimensional \(\mathbb{C}G\)-module. Then
    \[
        \langle \chi_V,\chi_V\rangle=1
    \]
    if and only if \(V\) is simple.
\end{corollary}

\begin{proof}
    Since \(\mathbb{C}G\) is semisimple, we may write
    \[
        V
        \cong
        S_1^{\oplus n_1}
        \oplus
        \cdots
        \oplus
        S_r^{\oplus n_r},
    \]
    where \(S_1,\dots,S_r\) are pairwise non-isomorphic simple
    \(\mathbb{C}G\)-modules and \(n_i\geq 0\). Then
    \[
        \chi_V
        =
        n_1\chi_{S_1}
        +
        \cdots
        +
        n_r\chi_{S_r}.
    \]
    By row orthogonality,
    \[
        \langle \chi_V,\chi_V\rangle
        =
        n_1^2+\cdots+n_r^2.
    \]
    Therefore
    \[
        \langle \chi_V,\chi_V\rangle=1
    \]
    if and only if exactly one of the \(n_i\)'s is equal to \(1\), and all the
    others are \(0\). This is equivalent to saying that \(V\) is simple.
\end{proof}

\begin{lemma}
    Let \(V\) and \(W\) be finite-dimensional \(\mathbb{C}G\)-modules. Suppose
    \[
        V
        \cong
        S_1^{\oplus n_1}
        \oplus
        \cdots
        \oplus
        S_r^{\oplus n_r},
    \]
    and
    \[
        W
        \cong
        S_1^{\oplus m_1}
        \oplus
        \cdots
        \oplus
        S_r^{\oplus m_r},
    \]
    where \(S_1,\dots,S_r\) are pairwise non-isomorphic simple
    \(\mathbb{C}G\)-modules and \(n_i,m_i\geq 0\). Then
    \[
        \langle \chi_V,\chi_W\rangle
        =
        n_1m_1+\cdots+n_rm_r.
    \]
\end{lemma}

\begin{proof}
    We have
    \[
        \chi_V
        =
        \sum_{i=1}^r n_i\chi_{S_i},
        \qquad
        \chi_W
        =
        \sum_{i=1}^r m_i\chi_{S_i}.
    \]
    Therefore, by bilinearity of the inner product and row orthogonality,
    \[
        \begin{aligned}
            \langle \chi_V,\chi_W\rangle
             & =
            \left\langle
            \sum_{i=1}^r n_i\chi_{S_i},
            \sum_{j=1}^r m_j\chi_{S_j}
            \right\rangle                        \\
             & =
            \sum_{i,j}n_i m_j
            \langle \chi_{S_i},\chi_{S_j}\rangle \\
             & =
            \sum_{i=1}^r n_i m_i.
        \end{aligned}
    \]
\end{proof}

\begin{corollary}
    Let \(V\) and \(W\) be finite-dimensional \(\mathbb{C}G\)-modules. Then
    \[
        \chi_V=\chi_W
    \]
    if and only if
    \[
        V\cong W.
    \]
\end{corollary}

\begin{proof}
    If \(V\cong W\), then clearly \(\chi_V=\chi_W\).

    Conversely, write
    \[
        V
        \cong
        S_1^{\oplus n_1}
        \oplus
        \cdots
        \oplus
        S_r^{\oplus n_r},
    \]
    and
    \[
        W
        \cong
        S_1^{\oplus m_1}
        \oplus
        \cdots
        \oplus
        S_r^{\oplus m_r}.
    \]
    If \(\chi_V=\chi_W\), then for each \(i\),
    \[
        n_i
        =
        \langle \chi_V,\chi_{S_i}\rangle
        =
        \langle \chi_W,\chi_{S_i}\rangle
        =
        m_i.
    \]
    Hence \(V\) and \(W\) have the same simple direct summands with the same
    multiplicities. Therefore
    \[
        V\cong W.
    \]
\end{proof}

\begin{corollary}
    Let \(V\) be a finite-dimensional \(\mathbb{C}G\)-module, and suppose
    \[
        V
        \cong
        S_1^{\oplus n_1}
        \oplus
        \cdots
        \oplus
        S_r^{\oplus n_r},
    \]
    where \(S_1,\dots,S_r\) are pairwise non-isomorphic simple
    \(\mathbb{C}G\)-modules. Then
    \[
        n_i
        =
        \langle \chi_V,\chi_{S_i}\rangle
    \]
    for every \(i\).
\end{corollary}

\begin{proof}
    Since
    \[
        \chi_V
        =
        \sum_{j=1}^r n_j\chi_{S_j},
    \]
    row orthogonality gives
    \[
        \langle \chi_V,\chi_{S_i}\rangle
        =
        \sum_{j=1}^r n_j\langle \chi_{S_j},\chi_{S_i}\rangle
        =
        n_i.
    \]
\end{proof}

Let
\[
    \operatorname{Irr}(G)=\{\chi_1,\dots,\chi_r\},
\]
and let
\[
    \operatorname{Cl}(G)=\{C_1,\dots,C_r\}
\]
be the set of conjugacy classes of \(G\). Choose a representative
\[
    g_j\in C_j
\]
for each \(j\). The character table of \(G\) is the matrix
\[
    X=(\chi_i(g_j))_{1\leq i,j\leq r}.
\]

\begin{theorem}[Column orthogonality relations]
    With the notation above, for \(1\leq j,\ell\leq r\), one has
    \[
        \sum_{i=1}^r
        \chi_i(g_j)\overline{\chi_i(g_\ell)}
        =
        \begin{cases}
            |C_G(g_j)|, & j=\ell,     \\
            0,          & j\neq \ell.
        \end{cases}
    \]
    Equivalently, for arbitrary \(g,h\in G\),
    \[
        \sum_{\chi\in \operatorname{Irr}(G)}
        \chi(g)\overline{\chi(h)}
        =
        \begin{cases}
            |C_G(g)|, & g \text{ and } h \text{ are conjugate in }G,     \\
            0,        & g \text{ and } h \text{ are not conjugate in }G.
        \end{cases}
    \]
\end{theorem}

\begin{proof}
    Let
    \[
        D
        =
        \frac{1}{|G|}
        \operatorname{diag}(|C_1|,\dots,|C_r|).
    \]
    The row orthogonality relations say that
    \[
        \langle \chi_i,\chi_\ell\rangle
        =
        \delta_{i\ell}.
    \]
    Writing the inner product as a sum over conjugacy classes gives
    \[
        \frac{1}{|G|}
        \sum_{j=1}^r
        |C_j|\chi_i(g_j)\overline{\chi_\ell(g_j)}
        =
        \delta_{i\ell}.
    \]
    In matrix form, this is
    \[
        X D \overline{X}^{\,t}=I_r.
    \]
    Hence \(X\) is invertible, and
    \[
        D\overline{X}^{\,t}=X^{-1}.
    \]
    Multiplying on the right by \(X\), we get
    \[
        D\overline{X}^{\,t}X=I_r.
    \]
    Therefore
    \[
        \overline{X}^{\,t}X
        =
        D^{-1}.
    \]
    But
    \[
        D^{-1}
        =
        \operatorname{diag}\left(
        \frac{|G|}{|C_1|},
        \dots,
        \frac{|G|}{|C_r|}
        \right).
    \]
    By the orbit-stabilizer formula for the conjugation action,
    \[
        |C_j|
        =
        \frac{|G|}{|C_G(g_j)|}.
    \]
    Hence
    \[
        D^{-1}
        =
        \operatorname{diag}\left(
        |C_G(g_1)|,
        \dots,
        |C_G(g_r)|
        \right).
    \]
    The \((j,\ell)\)-entry of \(\overline{X}^{\,t}X\) is
    \[
        \sum_{i=1}^r
        \overline{\chi_i(g_j)}\chi_i(g_\ell).
    \]
    Thus
    \[
        \sum_{i=1}^r
        \overline{\chi_i(g_j)}\chi_i(g_\ell)
        =
        \begin{cases}
            |C_G(g_j)|, & j=\ell,     \\
            0,          & j\neq \ell.
        \end{cases}
    \]
    Taking complex conjugates gives
    \[
        \sum_{i=1}^r
        \chi_i(g_j)\overline{\chi_i(g_\ell)}
        =
        \begin{cases}
            |C_G(g_j)|, & j=\ell,     \\
            0,          & j\neq \ell.
        \end{cases}
    \]

    Finally, if \(g,h\in G\) are arbitrary, then their character values depend only
    on their conjugacy classes. Therefore the formula above is equivalent to
    \[
        \sum_{\chi\in \operatorname{Irr}(G)}
        \chi(g)\overline{\chi(h)}
        =
        \begin{cases}
            |C_G(g)|, & g \text{ and } h \text{ are conjugate in }G,     \\
            0,        & g \text{ and } h \text{ are not conjugate in }G.
        \end{cases}
    \]
\end{proof}
